\section{Mathematical preliminaries}
\label{sec:prelim}

This section fixes notation and records four instruments used repeatedly below:
exponential waiting times and their races; continuous-time Markov chains; probability
generating functions; and first-step analysis. The treatment is deliberately short.
Full catalogues of discrete laws, Laplace--Stieltjes transforms, and simulation
algorithms are omitted; what is kept is the minimum needed for branching, quasi-stationarity,
and the absorption models of \cref{sec:moc}.

\subsection{Exponential clocks and competing risks}
\label{sec:exp-race}

In continuous time, events are organised by exponential clocks. A random variable
$X\sim\mathrm{Exp}(\lambda)$ with rate $\lambda>0$ has density
$f_X(x)=\lambda\ee^{-\lambda x}$ for $x\ge 0$, mean $1/\lambda$, and variance
$1/\lambda^2$. Among continuous distributions in elementary use it is memoryless: for
all $t_1,t_2>0$,
\begin{equation*}
\pr{X\ge t_1+t_2\mid X\ge t_1}=\pr{X\ge t_2}.
\end{equation*}
The residual waiting time does not depend on how long one has already waited. That
property is precisely what makes a rate-driven jump process Markovian.

When several independent exponential clocks run at once, two facts organise both
analysis and simulation. Let $X\sim\mathrm{Exp}(\lambda)$ and $Y\sim\mathrm{Exp}(\mu)$ be
independent. Then the first event time is again exponential,
\begin{equation}
\label{eq:minXY}
Z=\min\{X,Y\}\sim\mathrm{Exp}(\lambda+\mu),
\end{equation}
and the identity of the winner is a weighted choice independent of $Z$:
\begin{equation}
\label{eq:compXY}
\pr{X<Y}=\frac{\lambda}{\lambda+\mu},\qquad
\pr{Y<X}=\frac{\mu}{\lambda+\mu}.
\end{equation}
For $n$ independent clocks with rates $\lambda_1,\ldots,\lambda_n$ the same pattern
extends: the minimum is $\mathrm{Exp}(\sum_i\lambda_i)$, and clock $i$ wins with
probability $\lambda_i/\sum_j\lambda_j$. In the models of later chapters, birth, death,
absorption, and catastrophe are competing exponential channels of this kind.

\subsection{Continuous-time Markov chains}
\label{sec:ctmc}

A continuous-time Markov chain (CTMC) on a discrete state space $\Omega$ is a process
$(X(t))_{t\ge 0}$ for which the future depends on the past only through the present
state. Time-homogeneous transition probabilities are written
\begin{equation*}
p_{ij}(t)=\pr{X(t)=j\mid X(0)=i},\qquad i,j\in\Omega,
\end{equation*}
and models are specified by rates rather than by closed forms for $P(t)=(p_{ij}(t))$.
Differentiating short-time transitions produces the infinitesimal generator $Q=(q_{ij})$:
\begin{equation}
\label{eq:generator}
q_{ij}=\lim_{\Delta t\to 0^+}\frac{p_{ij}(\Delta t)}{\Delta t}\ (i\neq j),\qquad
q_{ii}=-\sum_{j\neq i}q_{ij},
\end{equation}
so that $p_{ij}(\Delta t)=\delta_{ij}+q_{ij}\Delta t+o(\Delta t)$. Off-diagonal entries
are non-negative; rows of $Q$ sum to zero.

Pathwise construction follows the exponential race of \cref{sec:exp-race}. While the
chain sits in a non-absorbing state $i$, it waits an exponential time of rate
$\Lambda_i=-q_{ii}$, then jumps to $j\neq i$ with probability $q_{ij}/\Lambda_i$. A state
$i$ is absorbing if $q_{ii}=0$. In linear birth--death models, population zero is
absorbing and represents extinction; in absorption models of \cref{sec:moc}, rupture or
uptake can introduce further absorbing compartments.

The forward Kolmogorov (master) equations for the occupation probabilities
$p_j(t)=\pr{X(t)=j}$ are the componentwise form of $\mathrm{d}P/\mathrm{d}t=QP$ (with
$P(0)=I$ in matrix notation). Probability flows into $j$ along incoming rates and out
along $\Lambda_j$. For infinite state spaces the master equations are rarely solved
state by state; generating functions compress them into a single partial differential
equation.

\subsection{Linear birth--death rates}
\label{sec:bd-linear-prelim}

The linear birth--death process on $\{0,1,2,\ldots\}$ is the standard population model
before catastrophe rates are added. From state $i\ge 1$, a birth moves to $i+1$ at rate
$i\lambda$ and a death moves to $i-1$ at rate $i\mu$, with $0$ absorbing. Each individual
reproduces and dies independently; the per-capita growth rate $\lambda-\mu$ decides
whether means explode or decay. Ultimate extinction from a single individual is certain
if $\lambda\le\mu$, and equals $\mu/\lambda$ if $\lambda>\mu$. Finite-time extinction
probabilities and the continuous-time quasi-stationary mean reappear in
\cref{sec:cts}. Birth--death--catastrophe dynamics --- an additional load-dependent
killing event that empties the population --- are defined briefly in \cref{sec:bdc} and
developed fully in later chapters.

\subsection{Probability generating functions}
\label{sec:pgf-prelim}

When counts are non-negative integers, the probability generating function packages an
entire distribution into a single analytic object. For a discrete random variable $X$
on $\{0,1,2,\ldots\}$,
\begin{equation}
\label{eq:pgf-def}
G_X(z)=\ex{z^X}=\sum_{k\ge 0}\pr{X=k}\,z^k,\qquad |z|\le 1.
\end{equation}
The operations used throughout are mechanical but decisive:
\begin{align*}
\pr{X=k}&=\frac1{k!}\,G_X^{(k)}(0),\\
G_X(1)&=1,\\
\ex{X}&=G_X'(1),\\
\mathrm{Var}(X)&=G_X''(1)+G_X'(1)-\bigl(G_X'(1)\bigr)^2.
\end{align*}
In particular $G_X(0)=\pr{X=0}$. When $X$ is a population size, that value is an
extinction probability --- the quantity small-population models are built to study.

For a non-negative integer-valued CTMC $X(t)$ with $X(0)=k$, the generating function
becomes a function of two variables,
\begin{equation}
\label{eq:pgf-process}
G_k(z,t)=\ex{z^{X(t)}\mid X(0)=k}=\sum_{i\ge 0}p_{ki}(t)\,z^i.
\end{equation}
Multiplying the Kolmogorov equations by $z^i$ and summing replaces an infinite system of
ordinary differential equations by a first-order partial differential equation for $G$.
That compression is the analytical centre of the absorption models in \cref{sec:moc}, and
of the birth--death--catastrophe analyses of later chapters.

When several types are tracked at once, the multivariate extension is
\begin{equation}
\label{eq:pgf-multi}
\phi_{\mathbf{X}}(z_1,\ldots,z_n)
=\sum_{k_1,\ldots,k_n\ge 0}
p(k_1,\ldots,k_n)\, z_1^{k_1}\cdots z_n^{k_n},
\end{equation}
with each $z_i\in[0,1]$. Multi-type generating functions are prepared ground for
two-type catastrophe models; they are not developed fully here.

\subsection{First-step analysis}
\label{sec:first-step}

When the question is not the full time course but a hitting probability or mean time,
first-step analysis conditions on the first exit from the current state and solves a
linear system. The same gesture appears as conditioning on which exponential clock won
the race, and as the Galton--Watson extinction fixed-point equation $s=\phi(s)$ of
\cref{sec:gw}: one step of randomness, then a self-consistent equation for the quantity
of interest.

With these instruments in place we turn to the discrete-generation branching process
that carries the first half of the chapter.
